\documentclass[12pt,a4paper]{article}%
\usepackage{amssymb}
\usepackage{amsfonts}
\usepackage{amsmath}
\usepackage{graphicx}
\usepackage[singlespacing]{setspace}
\usepackage{geometry}%
\setcounter{MaxMatrixCols}{30}
%TCIDATA{OutputFilter=latex2.dll}
%TCIDATA{Version=5.50.0.2960}
%TCIDATA{CSTFile=40 LaTeX article.cst}
%TCIDATA{Created=Saturday, April 12, 2014 14:08:00}
%TCIDATA{LastRevised=Friday, February 27, 2026 11:30:15}
%TCIDATA{<META NAME="GraphicsSave" CONTENT="32">}
%TCIDATA{<META NAME="SaveForMode" CONTENT="1">}
%TCIDATA{BibliographyScheme=Manual}
%TCIDATA{<META NAME="DocumentShell" CONTENT="Standard LaTeX\Standard LaTeX Article">}
%TCIDATA{Language=American English}
%TCIDATA{PageSetup=72,72,72,72,0}
%TCIDATA{ComputeDefs=
%$z_{t}$
%}
%BeginMSIPreambleData
\providecommand{\U}[1]{\protect\rule{.1in}{.1in}}
%EndMSIPreambleData
\newtheorem{theorem}{Theorem}
\newtheorem{acknowledgement}[theorem]{Acknowledgement}
\newtheorem{algorithm}[theorem]{Algorithm}
\newtheorem{axiom}[theorem]{Axiom}
\newtheorem{case}[theorem]{Case}
\newtheorem{claim}[theorem]{Claim}
\newtheorem{conclusion}[theorem]{Conclusion}
\newtheorem{condition}[theorem]{Condition}
\newtheorem{conjecture}[theorem]{Conjecture}
\newtheorem{corollary}[theorem]{Corollary}
\newtheorem{criterion}{Criterion}
\newtheorem{definition}{Definition}
\newtheorem{example}{Example}
\newtheorem{exercise}[theorem]{Exercise}
\newtheorem{lemma}[theorem]{Lemma}
\newtheorem{notation}[theorem]{Notation}
\newtheorem{problem}[theorem]{Problem}
\newtheorem{proposition}[theorem]{Proposition}
\newtheorem{remark}[theorem]{Remark}
\newtheorem{solution}[theorem]{Solution}
\newtheorem{summary}[theorem]{Summary}
\newenvironment{proof}[1][Proof]{\noindent\textbf{#1.} }{\ \rule{0.5em}{0.5em}}
\geometry{left=1in,right=1in,top=1in,bottom=1in}
\begin{document}
\textbf{HOMEWORK: MONETARY\ AND\ FINANCIAL\ ECONOMICS (Instructor: JB
Chatelain)}

The aim of this homework is to investigate (1) regime switching of
time-varying persistence (delay to revert to the mean equilibrium) and of (2)
regime-switching time-varying effects of feedback policy rules on policy
targets and (3) the identification issues of reverse causality due to negative
feedback leading to an opposite expected sign (for example, the price puzzle
in VAR models). There are examples of former homeworks on the EPI and some
former code with not exactly the same questions.

Answers can be written in french or english to be uploaded on the EPI on the
scheduled deadline day after feedback of your delegate. It is the
EPI\ "Monetary and Financial Macroeconomics" (access code for non-Paris1
registration is "monetary" or "Monetary"). You upload:

(1) Your full logbook in Word and also saved in pdf (with unlimited number of
pages) of artificial intelligence\ large language models (ChatGPT, Claude,
Deepseek, Mistral, Grok, Perplexity, Gemini,...) help, copying and pasting all
your prompts and answers while doing this homework. At the top, add 1 to
4\ pages summary at the very beginning highlighting what was mostly missing
from the AI answers. See more precisely question 1 on solving the price puzzle
with the help of artificial intelligence.

(2) Pdf file of your text answering question below, with at the end, a copy
and paste of your code and of its automatic translations using IA in the 3
other languages among the four languages: Python, R, Stata, SAS.\ 

(3) data set in Excel

(4) data set your software format, for example dta for Stata, for replication.

(5) code of the estimations run on the data set in your software format for
replication in STATA, R, Python \textbf{should be instantly run for
replication by the instructor using your code.}

The work is to be done individually. If you use monetary aggregates, you
should use their GROWTH rate. For other consumer or stock market prices
indices, you should use their GROWTH rates (inflation or returns) discrete or
continuous (differences of natural logarithm).

\textbf{Table 1: Suggestions for groups to be chosen, other countries (your
own country) and other variables welcome (monitoring exchange rate by the
Central Bank of emerging economies, for example):}%

\begin{tabular}
[t]{|c|c|c|c|c|}\hline
& Target $"\pi_{t}"$ & Instrument $"i_{t}"$ & Country & Frequency\\\hline
Afshar & Inflation CPI & Funds rate & US & Quarterly\\\hline
Battesti & Inflation\ CPI & Funds rate & Malta or Cyprus & Quarterly\\\hline
Bornier & Inflation\ CPI & Funds rate & France & Quarterly\\\hline
Esposito & Inflation CPI & Funds rate & Euro area & Quarterly\\\hline
Godo & Inflation\ CPI & Funds rate & Sweden & Quarterly\\\hline
Landini & Inflation\ CPI & Funds rate & Japan & Quarterly\\\hline
Lobo & Inflation CPI & Funds rate & Honduras & Quarterly\\\hline
Mitrofanoff & Inflation\ CPI & Funds rate & Germany & Quarterly\\\hline
Moncorge & Inflation\ CPI & Funds rate & Portugal & Quarterly\\\hline
Rahal & Inflation CPI & Funds rate & Poland & Quarterly\\\hline
Robertson & Inflation\ CPI & Funds rate & Serbia & Quarterly\\\hline
Uberti & Inflation\ CPI & Funds rate & UK & Quarterly\\\hline
Wei & Inflation CPI & Funds rate & China & Quarterly\\\hline
Wolff & Inflation CPI & Funds rate & Canada & Quarterly\\\hline
Yahiaoui & Inflation\ CPI & Funds rate & Italy & Quarterly\\\hline
Zhou & Inflation\ CPI & Funds rate & Spain & Quarterly\\\hline
17 & Private credit/GDP & Funds rate & UK & Quarterly\\\hline
18 & Public debt/GDP & Funds rate & US & Quarterly\\\hline
19 & Public debt/GDP & Funds rate & France & Quarterly\\\hline
20 & Public debt/GDP & Funds rate & Euro area & Quarterly\\\hline
21 & Public debt/GDP & Funds rate & UK & Quarterly\\\hline
22 & Public debt/GDP & Deficits,surplus/gdp & US & Quarterly\\\hline
23 & Public debt/GDP & Deficits,surplus/gdp & France & Quarterly\\\hline
24 & Public debt/GDP & Deficits,surplus/gdp & Euro area & Quarterly\\\hline
25 & Public debt/GDP & Deficits,surplus/gdp & UK & Quarterly\\\hline
\end{tabular}


You may propose two other time series involving feedback from another country
to the professor.

You can answer using the TEX file of the homework to begin with or copy it in
Word and use Word for your report and then transform it in pdf.

ADDITIONAL\ FILES on the EPI:

This text: Homework: in PDF, in TEX

EXAMPLES\ FOR\ HELP:

mavroeidis2010.dta , provides data in Stata format.

2018mavroeidisJBexamrolling09.do, Stata code for rolling windows regressions,
summary statistics and correlation.

2018 Violon HOMEWORK.pdf: old example of answers of the homework with some
DIFFERENT questions.

\section{Questions}

\begin{enumerate}
At the beginning, fill in this summary table once you are done with all questions:
\end{enumerate}%

\begin{tabular}
[t]{|c|c|c|}\hline
Country/frequency & USA & Quarterly\\\hline
Number of observ.: & Target $"\pi_{t}"$: & Instrument $"i_{t}"$:\\\hline
$T=160$ & Inflation CPI & Funds rate\\\hline
Beginning & 1981Q1 & 1981Q1\\\hline
End & 2020Q4 & 2020Q4\\\hline
AR(1) coeff & $\rho=0.97$ & $\rho=0.96$\\\hline
AR(1) of residuals & $0.50$ & $0.30$\\\hline
AR(2) lag1 coeff & $1.47$ & $1.30$\\\hline
AR(2) lag2 coeff & $-0.50$ & $-0.34$\\\hline
1st order SISO & $A=0.8$ & $F=0.5$\\\hline
& $B=0.1$ & $BF=0.05>0$\\\hline
Instrumental variables & Oil price (for example) & Romer shocks\\\hline
Indirect least squares & $B=0.3$ & $F=2$\\\hline
&  & $BF=0.6>0$\\\hline
\end{tabular}


\bigskip\textbf{PART 1: LLM\ logbook of all prompts and answers (an incomplete
logbook leads to reduction of the grade).}

\begin{enumerate}
\item Upload LLM logbook into a specific Word then saved as pdf files. Copy in
a separate Word file (uploaded as pdf on the EPI) with unlimited number of
pages: (1) At the top write your own final summary of what was finally useful
or missing using AI. (2) Then, you copy all the answers to all your prompts
from any LLM: Chatgpt, Deepseek, Claude, Mistral AI, Grok, Perplexity, Gemini.
Please copy and paste as many of your prompts and answers as possible for
code, text, translation, concepts, and so on. The aim is helping you to solve
the price puzzle in VAR and the statistical codes they propose and their
translation for Python, R, Stata and SAS.

\textbf{PART\ 2: Translation using LLM\ in order to have 3\ codes: Python,
Stata, R and database in format for instant running and replication (to be
done at the very end of your work).}

\item Translate your code so that you upload 3\ versions Python R and Stata of
code plus databases in required format for instant replication when running
the code.

\textbf{PART\ 3: Use LLM\ to address the VAR price puzzle due to the reverse
causality of monetary policy feedback rules (to be done first).}

\item Do a synthetic table of the 20 and more methods proposed for the prompt
"propose 20 detailed methods, ordered by their efficiency, to solve the price
puzzle using VAR models for each case with "why and how" precisions with a
related paper as a reference for each method [check the reference does exist]"
using recursively (copying the answer in another LLM asking for improvement)
for at least 4 LLM. Then ask the prompt "Sort your methods by categories for
example: (A) additional variables and more lags, (B) data transformation, (C)
Identification strategies, (D) different estimation methods of VAR (E)
others." Then ask prompt: "which theoretical mechanism and/or explain how and
why these methods may solve the price puzzle." Then ask prompts with correct
exact reference and download your preferred paper. Remark: there is a
meta-analysis by Havranek and others on the price puzzle.

(2) Then select one paper to replicate one estimation on Mendeley or Harvard
replication or on the website of the published paper or anywhere for the last
open question at the end of the homework. Explain why you select it. Pay
attention, you may need to download a few OTHER\ TIMES\ SERIES to apply their
method, check if these time series are available for your country and download
them at the same time than the CPI\ et the funds rate.

\textbf{PART 4: Your example}

\item Expand the data base for proxies of your policy target inflation and
your policy target for the longest period without spending too much time.
Provide the link to the database (FRED, ECB, BdF, BoE, IMF, WB and so on)
related to your time series. For example, do no use the CPI index but its
growth rate (or difference of log) which measure inflation. By default, you
may use quarterly frequency data instead of monthly, as in the last question,
you may expand variables adding quarterly output measures.

\item Check if there is seasonal adjustement or not: you may regress your time
series on seasonal dummies. Check if the scale of the time series is
comparable: for example, if monthly inflation is not "annualized" i.e.
multiplied by 12 to have an "annual" scale for this monthly data, whereas the
monthly Fed's funds rate is multiplied by 12. Then multiply monthly inflation
by 12, because in regressions, the ratio of standard errors of each variables
will be multiplied or divided by 12 because of these difference of scale. Then
the Taylor rule parameter may be multiplied by 12\ with monthly or by 4\ with
quarterly data.

\item Copy a plot of the two time series over time to check consistency with
answer 3. If you have done the correction, annualized funds rate should not be
too far from annualized inflation.

\item Report in a table univariate summary statistics: min max mean median
standard errors of policy target and policy instrument.

\item Estimate on the full period an AR(1) process and estimate the
auto-correlation of its residual (let's denote it: AR(1)-AR(1)). Remark: all
the equation taken from the cours are in deviation from the steady state mean,
you should keep everywhere intercepts for these means:%
\begin{align*}
\pi_{t+1} &  =\lambda_{\pi}\pi_{t}+u_{t+1}\text{ with }u_{t+1}=\rho_{\pi}%
u_{t}+\varepsilon_{t+1}\\
\text{ }i_{t+1} &  =\lambda_{i}i_{t}+u_{t+1}\text{ with }u_{t+1}=\rho_{i}%
u_{t}+\varepsilon_{t+1}%
\end{align*}


Additional: if you consider inflation also estimate the auto-regressive AR(1)
equation for output or output gap or economic growth (and conversely)%
\[
x_{t+1}=\lambda_{x}x_{t}+u_{t+1}\text{ with }u_{t+1}=\rho_{x}u_{t}%
+\varepsilon_{t+1}%
\]


\item Do the unit root tests for $\pi_{t}$ and $i_{t}$. Are they consistent
with your AR(1) estimations?

\item Estimate the AR(2) model. Check your estimates are close to the
following combination of parameters of the AR(1)-AR(1):
\begin{align*}
\pi_{t+1} &  =\left(  \lambda_{\pi}+\rho_{\pi}\right)  \pi_{t}-\lambda_{\pi
}\rho_{\pi}\pi_{t-1}+u_{t+1}\text{ with }u_{t+1}=\mu_{\pi}u_{t}+\varepsilon
_{t+1}\\
\text{ }i_{t+1} &  =\left(  \lambda_{i}+\rho_{i}\right)  i_{t}-\lambda_{i}%
\rho_{i}\pi_{t-1}+u_{t+1}\text{with }u_{t+1}=\mu_{i}u_{t}+\varepsilon_{t+1}%
\end{align*}


\item Compute the roots of the characteristic polynomial of your AR(2)
estimations. Are they real roots or complex conjugate? Are they within the
unit disk of complex numbers for stability. This is valid also for following
questions: Beware, some code may propose you the solution with a polynomial
where $Z=1/X$ with the inverse of the roots, so that for example $X^{2}$ in
the equation below is $1$ after simplication. As well for the roots of VAR,
nearly all your roots are inside the unit disk. If the code provides you all
of them outside the unit disk, you need to compute their inverse with their
proper value to be inside the unit disk, please, double check and think.
\[
X^{2}-\left(  \lambda_{\pi}+\rho_{\pi}\right)  X-\lambda_{\pi}\rho_{\pi}=0
\]


\item Draw the stability triangle for the AR(2) model with sum of the roots on
the horizontal axis and product of the roots on the vertical axis. Are the sum
and product within the stability triangle (so that the roots are within the
unit disk of complex numbers $\left\vert \rho_{\pi}\right\vert <1$ and
$\left\vert \lambda_{\pi}\right\vert <1$). Check whether the roots are complex
conjugate or two real roots in the stability triangle with respect to their
location of sum and product with respect to the quadratic curve corresponding
to the discriminant of the characteristic polynomial.

\item Check the partial autocorrelation order $p$ where the order $p+1$ leads
to a partial autocorrelation not statistically different from zero (i.e. test
the order of the AR(p) model).

\item For an AR(p), $p\geq2$, for measuring persistence, should you consider
the first auto-regressive parameter only or the sum of the auto-regressive
parameters? A\ PhD\ student (who has forgotten this exercise) comes to you and
states: look the auto-regressive parameter (the first one of the AR(2)) is
$1.4$, it is very much above one with a monthly or quarterly deterministic
trend with 40\% growth rate (but he does not discuss the second autoregressive
estimate which is $-0.4$). So he says, there is necessarily a unit root! What
do you answer to him?

\item In VAR\ estimation, are the residuals supposed to be white noise or
auto-regressive? What is the difference between weak white noise and strong
white noise? Which are the null hypothesis and the statistics of the Ljung-Box
(portmanteau) test?

\item DSGE models estimations add systematically auto-regressive shocks to
each equation. They give a particular "insightful" name of each of these
shocks (see Paul Romer critique of DSGE models). Half of these auto-regressive
parameters turn to be estimated to be close to unit root using "Bayesian"
estimation (see e.g. Smets and Wouters paper). Having not only persistent
residuals but also some residuals with unit root: is it expected to be valid
in VAR\ or VECM? Check various LLM\ answers on evidence or not that Bayesian
econometrics in DSGE\ models deny the importance of unit roots for time series
of shocks?

\item Explain the observational equivalence with AR(1)+AR(1) versus AR(2)? Is
this observational equivalence possible if the AR(2) process has two complex
conjugate roots?

\item Estimate with 10 years rolling windows the auto-correlation coeffcient
of AR(1) for the policy instrument and the policy target. Plot estimates over
time. Do you distinguish different periods for distinct values of persistence:
unit root or stationary process?

\item Let's investigate if we have at least three year periods of a regime
with trend inflation. Estimate three years rolling windows of the
deterministic trend equation for the policy target. Plot time-varying
estimates of the trend $g$ parameter as a function of time:
\[
Ln(\pi_{t+1})=a+gt+\varepsilon
\]


\item For unit root periods (auto-correlation close to one) of the policy
target, if you notice on the graph of the time series a rising trend subperiod
and estimate a determenistic trend on this subperiod: $Ln(\pi_{t+1}%
)=a+gt+\varepsilon$ ($g>0$). Do the same if you notice a subperiod with
declining trend: $Ln(\pi_{t+1})=a+gt+\varepsilon$ ($g<0$). Remark: $E_{0}%
\pi_{t+1}=\left(  1+g\right)  ^{t}\pi_{0}$ implies $Ln(E_{0}\pi_{t+1}%
)=Ln(\pi_{0})+tLn(1+g)\approx Ln(\pi_{0})+gt$ if $g$ small.

\item Test for unit roots for policy instrument and target (on distinct
periods if shifts of stationarity). Comment the results.

\item Provides tables of cross-correlation of the policy instrument $i_{t}$
with the policy target $\pi_{t+2}$, $\pi_{t+1}$, $\pi_{t}$,$\pi_{t-1}$%
,$\pi_{t-2}$. Plot 10 years rolling windows of these cross-correlations (if
possible). Comment the results.

\item Using 10 years rolling windows of the estimation of the first order
single input single output model (estimate also the auto-correlation of order
one of the residuals for each equation):%
\begin{align*}
\pi_{t+1} &  =A\pi_{t}+Bi_{t}+u_{t+1}\text{ with }u_{t+1}=\rho u_{t}%
+\varepsilon_{t+1}\\
\text{ }i_{t} &  =F\pi_{t}+u_{t+1}\text{ with }u_{t+1}=\rho u_{t}%
+\varepsilon_{t+1}%
\end{align*}


\item Using 10 years rolling windows, compute $A+BF$ and compare it to your
AR(1) estimate $\lambda_{\pi}$ of inflation:%
\begin{align*}
\pi_{t+1} &  =(A+BF)\pi_{t}+u_{t+1}\text{ with }u_{t+1}=\rho u_{t}%
+\varepsilon_{t+1}\\
\text{ }\pi_{t+1} &  =\lambda_{\pi}\pi_{t}+u_{t+1}\text{ with }u_{t+1}=\rho
u_{t}+\varepsilon_{t+1}%
\end{align*}


\item Plot $A$, $B$, $F$, $A+BF$ for rolling windows on the same graph. How is
this comparison related to the Lucas' critique ? Regarding the signs of
estimates of $B$ and $F$, do you conclude to positive feedback or negative
feedback ($0<A+BF<A$) taking into account the estimated values of the policy
rule parameter?

\item Redo the last 3\ questions with all variables in first differences (try
your luck).

\item Longest period single estimate. Here are the general formulas from
descriptive statistics to estimated parameters in the trivariate regression
case:%
\begin{align*}
\pi_{t+1} &  =\frac{r\left(  \pi_{t+1},\pi_{t}\right)  -r\left(  \pi
_{t+1},i_{t}\right)  r\left(  i,\pi_{t}\right)  }{1-r^{2}\left(  \pi_{t}%
,i_{t}\right)  }\frac{\sigma\left(  \pi_{t+1}\right)  }{\sigma\left(  \pi
_{t}\right)  }\pi_{t}+\frac{r\left(  \pi_{t+1},i_{t}\right)  -r\left(
\pi_{t+1},\pi_{t}\right)  r\left(  \pi_{t},i_{t}\right)  }{1-r^{2}\left(
\pi_{t},i_{t}\right)  }\frac{\sigma\left(  \pi_{t+1}\right)  }{\sigma\left(
i_{t}\right)  }i_{t}+u_{t+1}\\
i_{t} &  =r\left(  i_{t},\pi_{t}\right)  \frac{\sigma\left(  i_{t}\right)
}{\sigma\left(  \pi_{t}\right)  }\pi_{t}+v_{t}\text{ which decomposes:
}F=r\left(  i_{t},\pi_{t}\right)  \frac{\sigma\left(  i_{t}\right)  }%
{\sigma\left(  \pi_{t}\right)  }\\
\pi_{t+1} &  =r\left(  \pi_{t+1},\pi_{t}\right)  \frac{\sigma\left(  \pi
_{t+1}\right)  }{\sigma\left(  \pi_{t}\right)  }\pi_{t}+w_{t}\text{, which
decomposes: }\lambda_{\pi}=r\left(  i_{t},\pi_{t}\right)  \frac{\sigma\left(
i_{t}\right)  }{\sigma\left(  \pi_{t}\right)  }%
\end{align*}
Could you replace the numerical values of descriptive statistcs in each
formulas and comment from this decomposition and as well comment again on the
decomposition between $\lambda_{\pi}$ and $A+BF$ for the longest period single estimates?

\item Plot rolling windows estimates of simple correlation and standard errors
of $r(\pi_{t},i_{t})$ and the ratio of standard errors of the policy
instruments over the policy target ($\frac{\sigma_{\pi}}{\sigma_{i}}$) and of
the estimated $\widehat{F}_{\pi}$ parameter of the simple regression between
the policy instrument [you may compute this ratio and plot the three variables
over time using excel] over your extended period for windows of 10 (and/or 15
years). Comment and interpret the change of the feedback rule parameters over
time depending on the changes of the two times over time $r(\pi_{t},i_{t})$
and ($\frac{\sigma_{\pi}}{\sigma_{i}}$) in its decomposition (and in
particular when $F$ is larger than one):%
\[
F_{\pi}=r\left(  i_{t},\pi_{t}\right)  \frac{\sigma_{i}}{\sigma_{\pi}%
}>1\Leftrightarrow0<\frac{\sigma_{\pi}}{\sigma_{i}}<r\left(  i_{t},\pi
_{t}\right)  <1
\]


\item Indirect least square with simultaneous equations. It is very hard to
find the instruments with true zero restrictions in the structural model.
Indirect least square estimation: Solve the two structural parameters $B$ and
$F$ as rational functions of some of the four reduced form estimated
parameters $\left(  \beta_{\pi y},\beta_{\pi y},\beta_{\pi y},\beta_{\pi
y}\right)  $ of the reduced form system of equations $(S_{R})$:
\[
(S_{1}):\left\{
\begin{array}
[c]{c}%
\pi_{t}=Bi_{t}+Gy_{t}+0z_{t}+\varepsilon_{t}\\
i_{t}=F\pi_{t}+0y_{t}+Hz_{t}+\eta_{t}%
\end{array}
\right.  \Leftrightarrow(S_{R}):\left\{
\begin{array}
[c]{c}%
\pi_{t}=\beta_{\pi y}y_{t}+\beta_{\pi z}z_{t}+\varepsilon_{R,t}\\
i_{t}=\beta_{iy}y_{t}+\beta_{iz}z_{t}+\eta_{R,t}%
\end{array}
\right.
\]
You find:%
\[
\widehat{F}=\frac{\beta_{iy}}{\beta_{\pi y}}\text{ and }\widehat{B}%
=\frac{\beta_{\pi z}}{\beta_{iz}}%
\]
Upon your preferences, you seek and download variables $y_{t}$ and $z_{t}$ and
estimate $(S_{R})$ with identification restrictions in the system $(S_{1})$,
setting their parameter to zero in each of the two structural equations. You
recover estimates of $\widehat{B}$ and $\widehat{F}$. Your results have 98\%
chances of not being "succesful" for finding opposite signs on some periods
for $\widehat{B}$ and $\widehat{F}$. You have a large choice: congressional
budget office output gap, other measures of output gap, measures of output,
employment or unemployment, measures of core inflation, monetary aggregates,
indicators of the Fed balance sheet relative to quantitative easing, public
debt (divided by gdp or it growth rate), credit to the private sector, asset
prices, housing prices, exchange rate, price of oil, dummy variables for
changes of Fed chairman or US president, dummy variables from monetary or
fiscal policy shocks found in other studies, forecast errors with respect to
announcements taken from financial markets indicators or business surveys. If
you use other price indices, for inflation or returns, you need to compute the
first differences of natural logarithm or the growth rate of these time series.

\item Check that you do not have weak instruments in each equation of
$(S_{R})$ with $R^{2}>0.1$.

\item Try a rolling regression with 6 quarter lag between interest rat!e and
inflation in the hope of having a negative sign for $B$:%
\begin{align*}
\pi_{t+1} &  =A\pi_{t-1}+Bi_{t-6}+u_{t+1}\text{ with }u_{t+1}=\rho
u_{t}+\varepsilon_{t+1}\\
\text{ }i_{t} &  =F\pi_{t}+u_{t+1}\text{ with }u_{t+1}=\rho u_{t}%
+\varepsilon_{t+1}%
\end{align*}


\textbf{PART 5: Replication of one method by your preferred  reference paper
found after your LLM\ investigations}

\item Use one method of the paper you find to replicate on your data set to
run an estimation (you may add several varaibles) attempting to do a sign flip
including your policy target and your policy instrument (and possibly other
variables) in order to solve the price puzzle obtaining a sign flip $B<0$. It
is NOT\ a problem if you fail. If you run a VAR: (1) show the unit disk with
the roots of the VAR and comment if there are still roots close to unit roots
(2) Mention if the parameters of the lagged dependent variable remain close to
one (main diagonal), in this case, there remains very little variance to be
explained by other variables, even though those variables are statistically
significant with $T>200$.
\end{enumerate}

\textbf{APPENDIX: STATA\ CODE: }2018mavroeidisJBexamrolling09.do

*******************************************

* fyff federal funds rate and inflation

* Pre-Volcker: 1960:Q1 to 1979:Q2

* Greenspan's period: 1987:Q3 to 2006:Q1.

* generate: to be run only one time in a session

* CRTL D after selecting lines from top.

* capture log close

* log using mavroeidisJBexam.txt, text replace

*******************************************

set more off

version 10.1

clear all

set memory 30m

set linesize 90

cd D:%
%TCIMACRO{\TEXTsymbol{\backslash}}%
%BeginExpansion
$\backslash$%
%EndExpansion
Utilisateurs%
%TCIMACRO{\TEXTsymbol{\backslash}}%
%BeginExpansion
$\backslash$%
%EndExpansion
j.chatelain%
%TCIMACRO{\TEXTsymbol{\backslash}}%
%BeginExpansion
$\backslash$%
%EndExpansion
Documents%
%TCIMACRO{\TEXTsymbol{\backslash}}%
%BeginExpansion
$\backslash$%
%EndExpansion
\_Mes\_do1%
%TCIMACRO{\TEXTsymbol{\backslash}}%
%BeginExpansion
$\backslash$%
%EndExpansion
\_\_\_\_Rech%
%TCIMACRO{\TEXTsymbol{\backslash}}%
%BeginExpansion
$\backslash$%
%EndExpansion
\_\_\_\_Art-KJB%
%TCIMACRO{\TEXTsymbol{\backslash}}%
%BeginExpansion
$\backslash$%
%EndExpansion
\_0\_\_0closedloop\_Stata

** cd E:%
%TCIMACRO{\TEXTsymbol{\backslash}}%
%BeginExpansion
$\backslash$%
%EndExpansion


use mavroeidis2010

describe

gen newt = tq(1960q1) + t - 1

tsset newt, quarterly

describe

gen linflation = l.inflation

gen dinflation = inflation - linflation

gen lfyff = l.fyff

gen dfyff = fyff - lfyff

gen l2fyff = l.lfyff

gen lgapcbo = l.gapcbo

gen dgapcbo = gapcbo-lgapcbo

gen realrate=-1+((1+fyff)/(1+inflation))

gen lrealrate = lfyff - linflation

egen stdinflation = std(inflation)

egen stdfyff = std(fyff)

gen absstdinflation = abs(stdinflation)

gen absstdfyff = abs(stdfyff)

generate id1=(tt%
%TCIMACRO{\TEXTsymbol{>}}%
%BeginExpansion
$>$%
%EndExpansion
78)

generate id2=(tt%
%TCIMACRO{\TEXTsymbol{>}}%
%BeginExpansion
$>$%
%EndExpansion
88)

********** ROLLING WINDOW**********

* For any reason, one may need to begin again the code at the top

* for each rolling regression as variables may not be kept

* in memory. quarterly window(40) 10 years of 4 quarters

*****************************************

/*

* rolling 1st order SISO, policy instrument rule and autocorrelation of policy target

rolling, window(40) clear: nlsur
(inflation=\{binflation=0.56\}*linflation+\{bfyff=-0.04\}*lfyff+\{meaninflation\})
///

(fyff=\{Taylorinflation=-0.16\}*inflation+\{meanfyff\}) ///

(inflation=\{autocorrinflation=-0.16\}*linflation+\{b7\}) ///

if tt%
%TCIMACRO{\TEXTsymbol{>}}%
%BeginExpansion
$>$%
%EndExpansion
1

tsset start

list autocorrinflation binflation bfyff Taylorinflation meaninflation meanfyff

tsline autocorrinflation binflation bfyff Taylorinflation

* rolling correlation

rolling correlation=r(rho), window(40) clear: corr fyff inflation

tsset start

tsline correlation

* rolling correlation with mvcorr

ssc install mvcorr

mvcorr inflation fyff, window(40) gen(rho)

list rho

tsline rho

* rolling single standard deviation

rolling SD = r(sd), window(40) clear: summarize inflation

tsset start\qquad

tsline SD

* rolling single standard deviation

rolling SD = r(sd), window(40) clear: summarize fyff

tsset start

tsline SD

* rolling regression with output gap window(80)

rolling, window(60) clear: reg fyff lfyff linflation lgapcbo

tsset start

tsline \_b\_cons \_b\_lfyff \_b\_linflation \_b\_lgapcbo

* rolling Taylor rule with output gap window(80)

rolling, window(80) clear: reg fyff lfyff linflation lgapcbo

tsset start

tsline \_b\_cons \_b\_lfyff \_b\_linflation \_b\_lgapcbo

* rolling Taylor rule with output gap window(80)

reg inflation L(1/3).inflation, r

reg inflation L(1/1).inflation, r

rolling, window(80) clear: regress inflation L(1/1).inflation

tsset start

tsline \_b\_cons \_stat\_1

tsline \_stat\_1

*/

********************************************

* You may skip with /* */ the rolling part

* other tools

********************************************

****************************

* time series plot

****************************

tsline inflation fyff , tline(1979q3) tline(1982q1)

tsline inflation gapcbo , tline(1979q3) tline(1982q1)

********************************************

* univariate statistics

********************************************

sort id1 tt

by id1: summarize inflation fyff realrate

by id1: tabstat inflation fyff realrate, ///

stat (count p50 mean sd min max skew kurtosis) col(stat)

******************************************************

* Taylor principle, serial correlation of residuals

* bivariate plots

********************************************************

sort id1 tt

by id1: pwcorr inflation fyff gapcbo

* SORT by time-series to obtain lag: tsset t,q

* conflict with sorting by structural break by id1

tsset newt, quarterly

* auto-correlation of residuals

* of policy instrument rule on inflation

reg fyff inflation if id1==0

predict resid0 if id1==0, residual

tabstat resid0 if id1==0, stat (count mean sd variance)

reg resid0 l.resid0 if id1==0

pwcorr resid0 l.resid0 if id1==0

estat durbinalt, small lags(1/2)

estat bgodfrey, small lags(1/2)

arima resid0, arima(1,0,0)

gen lresid0 = l.resid0

reg fyff inflation if id1==1

predict resid1 if id1==1, residual

tabstat resid1 if id1==1, stat (count mean sd variance)

reg resid1 l.resid1 if id1==1

pwcorr resid1 l.resid1 if id1==1

estat durbinalt, small lags(1/2)

estat bgodfrey, small lags(1/2)

gen lresid1 = l.resid1

* PROC GPLOT BIVARIATE LINEAR QUADRATIC LINE PROC LOWESS

graph twoway (scatter fyff inflation if id1==0, msize(large) msymbol(o)) ///

(lfit fyff inflation if id1==0, clstyle(p3) lwidth(vvvthick)), ///

plotregion(style(none)) ///

title("Federal Funds Rate and inflation 1960q1-1979q2") ///

xtitle("Inflation", size(medlarge)) xscale(titlegap(*5)) ///

ytitle("Federal funds rate", size(medlarge)) yscale(titlegap(*5)) ///

legend(pos(4) ring(0) col(1)) legend(size(small)) ///

legend(label(1 "Actual Data") label(2 "Linear fit") label(3 "Quadratic fit")
label(4 "Lowess"))

graph twoway (scatter resid0 lresid0 if id1==0, msize(large) msymbol(o)) ///

(lfit resid0 lresid0 if id1==0, clstyle(p3) lwidth(vvvthick)), ///

plotregion(style(none)) ///

title("Serial correlation of residuals 1960q1-1979q2") ///

xtitle("lag1 residuals of Federal Funds Rate on Inflation", size(medlarge))
xscale(titlegap(*5)) ///

ytitle("residuals of Federal Funds Rate on Inflation", size(medlarge))
yscale(titlegap(*5)) ///

legend(pos(4) ring(0) col(1)) legend(size(small)) ///

legend(label(1 "Actual Data") label(2 "Linear fit") label(3 "Quadratic fit")
label(4 "Lowess"))

graph twoway (scatter inflation linflation if id1==0, msize(large) msymbol(o)
ysc(r(0 8.2)) ) ///

(lfit inflation linflation if id1==0, clstyle(p3) lwidth(vvvthick)), ///

plotregion(style(none)) ///

title("Serial correlation of inflation id1==0") ///

xtitle("lag1 inflation", size(medlarge)) xscale(titlegap(*5)) ///

ytitle("inflation", size(medlarge)) yscale(titlegap(*5)) ///

legend(pos(4) ring(0) col(1)) legend(size(small)) ///

legend(label(1 "Actual Data") label(2 "Linear fit") label(3 "Quadratic fit")
label(4 "Lowess"))

graph twoway (scatter fyff lfyff if id1==0, msize(large) msymbol(o)) ///

(lfit fyff lfyff if id1==0, clstyle(p3) lwidth(vvvthick)), ///

plotregion(style(none)) ///

title("Serial correlation of Federal Funds Rate if id1==0") ///

xtitle("lag1 output gap", size(medlarge)) xscale(titlegap(*5)) ///

ytitle("output gap", size(medlarge)) yscale(titlegap(*5)) ///

legend(pos(4) ring(0) col(1)) legend(size(small)) ///

legend(label(1 "Actual Data") label(2 "Linear fit") label(3 "Quadratic fit")
label(4 "Lowess"))

* PROC GPLOT BIVARIATE LINEAR QUADRATIC LINE PROC LOWESS id1==1

graph twoway (scatter fyff inflation if id1==1, msize(large) msymbol(o)) ///

(lfit fyff inflation if id1==1, clstyle(p3) lwidth(vvvthick)), ///

plotregion(style(none)) ///

title("Federal Funds Rate and inflation if id1==1") ///

xtitle("Inflation", size(medlarge)) xscale(titlegap(*5)) ///

ytitle("Federal funds rate", size(medlarge)) yscale(titlegap(*5)) ///

legend(pos(4) ring(0) col(1)) legend(size(small)) ///

legend(label(1 "Actual Data") label(2 "Linear fit") )

graph twoway (scatter resid1 lresid1 if id1==1, msize(large) msymbol(o)) ///

(lfit resid1 lresid1 if id1==1, clstyle(p3) lwidth(vvvthick)), ///

plotregion(style(none)) ///

title("Serial correlation of residuals if id1==1") ///

xtitle("lag1 residuals of Federal Funds Rate on Inflation", size(medlarge))
xscale(titlegap(*5)) ///

ytitle("residuals of Federal Funds Rate on Inflation", size(medlarge))
yscale(titlegap(*5)) ///

legend(pos(4) ring(0) col(1)) legend(size(small)) ///

legend(label(1 "Actual Data") label(2 "Linear fit") )

graph twoway (scatter inflation linflation if id1==1, msize(large) msymbol(o)
ysc(r(0 8.2)) ) ///

(lfit inflation linflation if id1==1, clstyle(p3) lwidth(vvvthick)), ///

plotregion(style(none)) ///

title("Serial correlation of inflation id1==1") ///

xtitle("lag1 inflation", size(medlarge)) xscale(titlegap(*5)) ///

ytitle("inflation", size(medlarge)) yscale(titlegap(*5)) ///

legend(pos(4) ring(0) col(1)) legend(size(small)) ///

legend(label(1 "Actual Data") label(2 "Linear fit") label(3 "Quadratic fit")
label(4 "Lowess"))

graph twoway (scatter fyff lfyff if id1==1, msize(large) msymbol(o)) ///

(lfit fyff lfyff if id1==1, clstyle(p3) lwidth(vvvthick)), ///

plotregion(style(none)) ///

title("Serial correlation of Federal Funds Rate if id1==1") ///

xtitle("lag1 output gap", size(medlarge)) xscale(titlegap(*5)) ///

ytitle("output gap", size(medlarge)) yscale(titlegap(*5)) ///

legend(pos(4) ring(0) col(1)) legend(size(small)) ///

legend(label(1 "Actual Data") label(2 "Linear fit") label(3 "Quadratic fit")
label(4 "Lowess"))

***************************************

* AR(1)

***************************************

sort id1

by id1: pwcorr inflation linflation

by id1: reg inflation linflation

by id1: pwcorr fyff lfyff

by id1: reg fyff lfyff

tsset newt, quarterly

corrgram inflation if id2==0, lag(8)

ac inflation if id2==0, lag(8)

corrgram inflation if id2==1, lag(8)

ac inflation if id2==1, lag(8)

corrgram fyff if id2==0, lag(4)

ac fyff if id2==0, lag(4)

corrgram fyff if id2==1, lag(4)

ac fyff if id2==1, lag(4)

arima inflation if id1==0, arima(1,0,0)

reg inflation linflation if id1==0

predict arinfl0 if id1==0, residual

tabstat arinfl0 if id1==0, stat (count mean sd variance)

pwcorr arinfl0 l.arinfl0 if id1==0

reg arinfl0 l.arinfl0 if id1==0

arima arinfl0, arima(1,0,0)

arima inflation if id1==1, arima(1,0,0)

reg inflation linflation if id1==1

predict arinfl1 if id1==1, residual

tabstat arinfl1 if id1==1, stat (count mean sd variance)

pwcorr arinfl1 l.arinfl1 if id1==1

reg arinfl1 l.arinfl1 if id1==1

arima arinfl1, arima(1,0,0)

arima fyff if id1==0, arima(1,0,0)

reg fyff lfyff if id1==0

predict arfyff0 if id1==0, residual

tabstat arfyff0 if id1==0, stat (count mean sd variance)

pwcorr arfyff0 l.arfyff0 if id1==0

reg arfyff0 l.arfyff0 if id1==0

arima arfyff0, arima(1,0,0)

arima fyff if id1==1, arima(1,0,0)

reg fyff lfyff if id1==1

predict arfyff1 if id1==1, residual

tabstat arfyff1 if id1==1, stat (count mean sd variance)

pwcorr arfyff1 l.arfyff1 if id1==1

reg arfyff1 l.arfyff1 if id1==1

arima arfyff1, arima(1,0,0)

********************

* unit root tests

************************

dfuller inflation if id1==0, notrend lag(1)

dfuller fyff if id1==0, notrend lag(1)

dfuller inflation if id1==1, notrend lag(1)

dfuller fyff if id1==1, notrend lag(1)

pperron inflation if id1==0, lag(1)

pperron fyff if id1==0, lag(1)

pperron inflation if id1==1, lag(1)

pperron fyff if id1==1, lag(1)

dfgls inflation if id1==0, notrend maxlag(1)

dfgls fyff if id1==0, notrend maxlag(1)

dfgls inflation if id1==1, notrend maxlag(1)

dfgls fyff if id1==1, notrend maxlag(1)

********************************************

* unconstrained VAR for comparison

********************************************

* VAR non contraint with constants

sort id2

by id2: var inflation fyff, lag(1 2)

varstable, graph

vargranger

varlmar, mlag(3)

sort id1

var inflation fyff if id1==0, lag(1)

varstable, graph

vargranger

varlmar, mlag(3)

predict residinf00, equation(\#1) residuals

predict residfyff00, equation(\#2) residuals

summarize residinf00 residfyff00 if id1==0

sort id1

reg residinf00 l.residinf00 linflation lfyff if id1==0

reg residinf00 l.residinf00 if id1==0

pwcorr residinf00 l.residinf00 if id1==0

reg residfyff00 l.residfyff00 linflation lfyff if id1==0

reg residfyff00 l.residfyff00 if id1==0

pwcorr residfyff00 l.residfyff00 if id1==0

* Verifie NLSUR=VAR with constants

nlsur ///

(inflation=\{b1=0.56\}*linflation+\{b2=-0.04\}*lfyff+\{b3\}) ///

(fyff=\{b4=-0.16\}*linflation+\{b5=0.91\}*lfyff+\{b6\}) ///

if tt%
%TCIMACRO{\TEXTsymbol{>}}%
%BeginExpansion
$>$%
%EndExpansion
1 \& id1==0

*****************************

* unconstrained VAR(1)

*****************************

var inflation fyff if id1==0, lag(1)

varstable, graph

vargranger

varlmar, mlag(1)

var inflation fyff if id1==1, lag(1)

varstable, graph

vargranger

varlmar, mlag(1)

var inflation fyff if id1==0, lag(1)

predict evarinfl0 if id1==0, residuals equation(inflation)

predict evarfyff0 if id1==0, residuals equation(fyff)

tabstat evarinfl0 evarfyff0 if id1==0, stat (count mean sd variance)

pwcorr evarinfl0 l.evarinfl0 evarfyff0 l.evarfyff0 if id1==0

reg evarinfl0 l.evarinfl0 if id1==0

reg evarfyff0 l.evarfyff0 if id1==0

var inflation fyff if id1==1, lag(1)

predict evarinfl1 if id1==1, residuals equation(inflation)

predict evarfyff1 if id1==1, residuals equation(fyff)

tabstat evarinfl1 evarfyff1 if id1==1, stat (count mean sd variance)

pwcorr evarinfl1 l.evarinfl1 evarfyff1 l.evarfyff1 if id1==1

reg evarinfl1 l.evarinfl1 if id1==1

reg evarfyff1 l.evarfyff1 if id1==1

* IRF FIGURES VAR(1)

varbasic inflation fyff if id1==0, lag(1)

irf graph fevd, lstep(1)

varbasic inflation fyff if id1==1, lag(1)

irf graph fevd, lstep(1)

**********************************************************

* cross correlogram inflation federal funds rate

**********************************************************

sort id1

by id1: pwcorr linflation fyff

by id1: pwcorr inflation fyff

by id1: pwcorr inflation lfyff

tsset newt, quarterly

xcorr inflation fyff if id1==0, table lags(1)

xcorr inflation fyff if id1==1, table lags(1)

xcorr inflation fyff if id1==0, lags(4)

xcorr inflation fyff if id1==1, lags(4)

********************************************

* Estimate A+BF, A, B and F over two subperiods

********************************************

* Verifie VAR non contraint with constants

pwcorr inflation gapcbo fyff realrate if id1==0

var inflation fyff gapcbo if id1==0, lag(1)

var gapcbo if id1==0, lag(1)

sort id2

by id2: pwcorr inflation linflation gapcbo lgapcbo fyff lfyff realrate lrealrate

sort id1

by id1: pwcorr inflation linflation gapcbo lgapcbo fyff lfyff realrate lrealrate

* Estimate A+BF, A, B and F over two subperiods

sort id1

by id1: nlsur ///

(inflation=\{autoinfl=0.56\}*linflation+\{cinflAR1\}) ///

(inflation=\{A=0.56\}*linflation+\{B=-0.04\}*lfyff+\{cinfl\}) ///

(fyff=\{F=-0.16\}*inflation+\{cfyff\}) ///

if tt%
%TCIMACRO{\TEXTsymbol{>}}%
%BeginExpansion
$>$%
%EndExpansion
1

sort id2

by id2: nlsur ///

(inflation=\{autoinfl=0.56\}*linflation+\{cinflAR1\}) ///

(inflation=\{A=0.56\}*linflation+\{B=-0.04\}*lfyff+\{cinfl\}) ///

(fyff=\{F=-0.16\}*inflation+\{cfyff\}) ///

if tt%
%TCIMACRO{\TEXTsymbol{>}}%
%BeginExpansion
$>$%
%EndExpansion
1

sort id2

by id2: pwcorr inflation linflation fyff lfyff gapcbo lgapcbo

sort id2

by id2: ///

var inflation fyff, lag(1)

varstable

* add output gap in transmission mechanism of inflation: not relevant

sort id2

by id2: nlsur ///

(inflation=\{autoinfl=0.56\}*linflation+\{cinflAR1\}) ///

(inflation=\{Ainfl\}*linflation+\{Bfyff\}*lfyff+\{Bgap=-0.04\}*lgapcbo+\{cinfla\})
///

(fyff=\{Finlf=-0.16\}*inflation+\{Fgap\}*gapcbo+\{cfyff\}) ///

(gapcbo=\{autogap\}*lgapcbo+\{cgapAR1\}) ///

(gapcbo=\{Agap\}*lgapcbo+\{Bfyffgap\}*lfyff+\{Binflgap=-0.04\}*linflation+\{cgap\})
///

if tt%
%TCIMACRO{\TEXTsymbol{>}}%
%BeginExpansion
$>$%
%EndExpansion
1

* real rate in transmission: negative sign Volcker

sort id2

by id2: nlsur
(inflation=\{binfl\}*linflation+\{brealrate\}*realrate+\{bgap=-0.04\}*gapcbo+\{cinfla\})
///

(fyff=\{b4=-0.16\}*inflation+\{b6\}) if tt%
%TCIMACRO{\TEXTsymbol{>}}%
%BeginExpansion
$>$%
%EndExpansion
1

sort id1

by id1: nlsur
(inflation=\{binfl\}*linflation+\{brealrate\}*realrate+\{bgap=-0.04\}*gapcbo+\{cinfla\})
///

(fyff=\{b4=-0.16\}*inflation+\{b6\}) if tt%
%TCIMACRO{\TEXTsymbol{>}}%
%BeginExpansion
$>$%
%EndExpansion
1

*********************************************************

* TAYLOR RULE with 2 lags auto-correlation of residuals

**********************************************************

reg fyff lfyff l2fyff linflation lgapcbo if id1==0

predict residinf20, residuals

reg residinf20 l.residinf20 linflation lgapcbo if id1==0

reg residinf20 l.residinf20 if id1==0

pwcorr residinf20 l.residinf20 if id1==0

estat durbinalt, small lags(1/2)

estat bgodfrey, small lags(1/2)

reg fyff linflation lgapcbo if id1==0

predict residinf201, residuals

reg residinf201 l.residinf201 linflation lgapcbo if id1==0

reg residinf201 l.residinf201 if id1==0

pwcorr residinf201 l.residinf201 if id1==0

estat durbinalt, small lags(1/2)

estat bgodfrey, small lags(1/2)

reg fyff lfyff l2fyff if id1==0

predict residinf202, residuals

reg residinf202 l.residinf202 linflation lgapcbo if id1==0

reg residinf202 l.residinf202 if id1==0

pwcorr residinf202 l.residinf202 if id1==0

estat durbinalt, small lags(1/2)

estat bgodfrey, small lags(1/2)

reg fyff lfyff linflation if id1==0

predict residinf203, residuals

reg residinf203 l.residinf203 linflation lgapcbo if id1==0

reg residinf203 l.residinf203 if id1==0

pwcorr residinf203 l.residinf203 if id1==0

estat durbinalt, small lags(1/2)

estat bgodfrey, small lags(1/2)

reg fyff lfyff lgapcbo if id1==0

predict residinf204, residuals

reg residinf204 l.residinf204 linflation lgapcbo if id1==0

reg residinf204 l.residinf204 if id1==0

pwcorr residinf204 l.residinf204 if id1==0

estat durbinalt, small lags(1/2)

estat bgodfrey, small lags(1/2)

* TAYLOR RULE

*by id1: reg fyff inflation gapcbo

*by id1: reg fyff lfyff inflation gapcbo

*by id1: nlsur (fyff=\{l\}*lfyff+(1-\{l\})*\{Fpi\}*inflation) if tt%
%TCIMACRO{\TEXTsymbol{>}}%
%BeginExpansion
$>$%
%EndExpansion
1

reg fyff inflation gapcbo if id1==0

predict resid01 if id1==0, residual

tabstat resid01 if id1==0, stat (count mean sd variance)

pwcorr resid01 l.resid01 if id1==0

reg resid01 l.resid01 if id1==0

arima resid01, arima(1,0,0)

reg fyff lfyff l2fyff inflation gapcbo if id1==0

estat durbinalt, small lags(1/2)

estat bgodfrey, small lagss(1/2)
\end{document}